\paragraph*{04.03.02.01 Die Grundlagen des Projektmanagements und deren Einführung kennen}
\kcilabel{para:04.03.02.01_GrundlagenProjektmanagement}{04.03.02.01}
\label{para:04.03.02.01_GrundlagenProjektmanagement}

%%%% Task
\par Als \gls{wsl} etablierte ich belastbare Governance-Strukturen für einheitliche Qualität und Termintreue bei 7 Zulieferern und 400 Objekten trotz geringem \gls{pmr}.

%%% Action
\par Ich definierte ein \gls{rasci}-basiertes Governance-Modell mit standardisierten Prozessen, Quality-Gates und Dokumentationsstandards. Implementierte wöchentliche Gremien (\gls{ais}-\gls{jf}, \gls{r2h}-\gls{jf}, \gls{f2h}-Schnittstellen-\gls{jf}) und steuerte den Übergang von Ad-hoc zu Meilenstein-basiertem Management.

%Result
\par Der \gls{pmr} stieg signifikant mit klaren Verantwortlichkeiten. Die Governance ermöglichte kontrollierte Führung der Implementierung durch Build, Test und Transport und reduzierte Reibungsverluste. Das \gls{wsr} blieb lieferfähig und unterstützte die termingerechte Migration.

%Reflect
\par Die frühe Etablierung von Governance war entscheidend für Transparenz und Koordination. In zukünftigen ähnlich komplexen Szenarien werde ich noch stärker auf regelmäßige Retrospektiven setzen, um Governance-Anpassungen agile vornehmen zu können.